% GENERATED FILE. Edit canonical lessons, then run npm run generate:books.
% canonical-source-hash: fnv1a64:de59f0470a54b3c3
% canonical-lessons: BN-C396-majhi, BN-C396-dar, BN-C396-nongor, BN-C396-chhoi, BN-C396-baoya

\chapter{Boats and Boatmen}
\label{ch:bn-boats-and-boatmen}
\hlchaptermodality{396}

\begin{chapteropening}
\textbf{By the end of this chapter:} \emph{I can talk about a boatman, his oars and anchor, and rowing across.}

Complete the last lesson of chapter 396: I can talk about a boatman, his oars and anchor, and rowing across.
\end{chapteropening}

\section[\texorpdfstring{\bn{মাঝি} (mājhi)}{মাঝি (mājhi)}]{\bn{মাঝি} (mājhi) — a boatman}

\label{lesson:BN-C396-majhi}



\begin{quote}
Before the new one: say the Bengali for abroad, then the Bengali for to make a crossing.
\end{quote}

\subsection*{You'll want to know: \bn{মাঝি}}
\textbf{\bn{মাঝি}} — \emph{mājhi} — \textquotedblleft{}a boatman\textquotedblright{}.

\begin{quote}
\bn{বুড়ো} \bn{মাঝি} \bn{নৌকায়} \bn{বসে} \bn{গান} \bn{গাইছিল।}
\end{quote}

\begin{quote}
\emph{buṛo mājhi noukāy bôse gān gāichhilô.}
\end{quote}

\begin{quote}
The old boatman sat in the boat, singing.
\end{quote}

The \bn{মাঝি}'s river songs are a whole branch of Bengali folk music.

\begin{cousinweb}
\textbf{\bn{মাঝি}} (\emph{mājhi}) is inherited: CDIAL lists Bengali \emph{māji}, \textquotedblleft{}a steersman\textquotedblright{}, under the reconstructed \emph{majjhika}, \textquotedblleft{}a boatman\textquotedblright{} (CDIAL 9714). It is not filed under \emph{mádhya}, middle, the source of \bn{মাঝ}.
\end{cousinweb}

\begin{grammarlens}[title={What you've built}]
One more word for getting around.
\end{grammarlens}

\subsection*{Your turn}
\begin{itemize}
\raggedright
  \item \emph{Say it:} \emph{mājhi}
  \item \emph{Say it:} \emph{mājhi}, once more
  \item \emph{Say it:} \emph{buṛo mājhi noukāy bôse gān gāichhilô.} — The old boatman sat in the boat, singing.
  \item \emph{Recall:} say the Bengali for abroad, then the Bengali for to make a crossing, then say \emph{mājhi} again
\end{itemize}

\begin{tcolorbox}[breakable,colback=teal!4,colframe=teal!35!black,title={Before you move on}]
What does \bn{মাঝি} mean? (\textquotedblleft{}A boatman\textquotedblright{}.) Say it once more.
\end{tcolorbox}
\section[\texorpdfstring{\bn{দাঁড়} (dā̃ṛ)}{দাঁড় (dā̃ṛ)}]{\bn{দাঁড়} (dā̃ṛ) — an oar}

\label{lesson:BN-C396-dar}



\begin{quote}
Before the new one: say the Bengali for to make a crossing, then the Bengali for a boatman.
\end{quote}

\subsection*{You'll want to know: \bn{দাঁড়}}
\textbf{\bn{দাঁড়}} — \emph{dā̃ṛ} — \textquotedblleft{}an oar\textquotedblright{}.

\begin{quote}
\bn{মাঝি} \bn{দুই} \bn{হাতে} \bn{দাঁড়} \bn{টানল।}
\end{quote}

\begin{quote}
\emph{mājhi dui hāte dā̃ṛ ṭānlô.}
\end{quote}

\begin{quote}
The boatman pulled at the oar with both hands.
\end{quote}

Not \bn{দাঁড়ানো}, to stand: the same letters, a different word.

\begin{cousinweb}
\textbf{\bn{দাঁড়}} (\emph{dā̃ṛ}) is inherited: CDIAL lists Bengali \emph{dā̃ṛ}, \textquotedblleft{}an oar\textquotedblright{}, under Sanskrit \emph{daṇḍá}, \textquotedblleft{}a stick, a staff\textquotedblright{} (CDIAL 6128): an oar is a stick with a blade.
\end{cousinweb}

\begin{grammarlens}[title={What you've built}]
One more word for getting around.
\end{grammarlens}

\subsection*{Your turn}
\begin{itemize}
\raggedright
  \item \emph{Say it:} \emph{dā̃ṛ}
  \item \emph{Say it:} \emph{dā̃ṛ}, once more
  \item \emph{Say it:} \emph{mājhi dui hāte dā̃ṛ ṭānlô.} — The boatman pulled at the oar with both hands.
  \item \emph{Recall:} say the Bengali for to make a crossing, then the Bengali for a boatman, then say \emph{dā̃ṛ} again
\end{itemize}

\begin{tcolorbox}[breakable,colback=teal!4,colframe=teal!35!black,title={Before you move on}]
What does \bn{দাঁড়} mean? (\textquotedblleft{}An oar\textquotedblright{}.) Say it once more.
\end{tcolorbox}
\section[\texorpdfstring{\bn{নোঙর} (nongôr)}{নোঙর (nongôr)}]{\bn{নোঙর} (nongôr) — an anchor}

\label{lesson:BN-C396-nongor}



\begin{quote}
Before the new one: say the Bengali for a boatman, then the Bengali for an oar.
\end{quote}

\subsection*{You'll want to know: \bn{নোঙর}}
\textbf{\bn{নোঙর}} — \emph{nongôr} — \textquotedblleft{}an anchor\textquotedblright{}.

\begin{quote}
\bn{বন্দরের} \bn{কাছে} \bn{জাহাজটা} \bn{নোঙর} \bn{ফেলল।}
\end{quote}

\begin{quote}
\emph{bôndôrer kāchhe jāhājṭā nongôr phellô.}
\end{quote}

\begin{quote}
The ship dropped anchor near the port.
\end{quote}

\bn{নোঙর} \bn{ফেলা}, to drop anchor; \bn{নোঙর} \bn{তোলা}, to weigh anchor.

\begin{cousinweb}
\textbf{\bn{নোঙর}} (\emph{nongôr}) is usually traced to Persian \emph{langar}, \textquotedblleft{}an anchor\textquotedblright{}; Bengali turned the first l into n.
\end{cousinweb}

\begin{grammarlens}[title={What you've built}]
One more word for getting around.
\end{grammarlens}

\subsection*{Your turn}
\begin{itemize}
\raggedright
  \item \emph{Say it:} \emph{nongôr}
  \item \emph{Say it:} \emph{nongôr}, once more
  \item \emph{Say it:} \emph{bôndôrer kāchhe jāhājṭā nongôr phellô.} — The ship dropped anchor near the port.
  \item \emph{Recall:} say the Bengali for a boatman, then the Bengali for an oar, then say \emph{nongôr} again
\end{itemize}

\begin{tcolorbox}[breakable,colback=teal!4,colframe=teal!35!black,title={Before you move on}]
What does \bn{নোঙর} mean? (\textquotedblleft{}An anchor\textquotedblright{}.) Say it once more.
\end{tcolorbox}
\section[\texorpdfstring{\bn{ছই} (chhôi)}{ছই (chhôi)}]{\bn{ছই} (chhôi) — the hood of a country boat}

\label{lesson:BN-C396-chhoi}



\begin{quote}
Before the new one: say the Bengali for an oar, then the Bengali for an anchor.
\end{quote}

\subsection*{You'll want to know: \bn{ছই}}
\textbf{\bn{ছই}} — \emph{chhôi} — \textquotedblleft{}the hood of a country boat\textquotedblright{}.

\begin{quote}
\bn{রোদে} \bn{আমরা} \bn{নৌকার} \bn{ছইয়ের} \bn{নিচে} \bn{বসলাম।}
\end{quote}

\begin{quote}
\emph{rode āmrā noukār chhôiyer niche bôslām.}
\end{quote}

\begin{quote}
In the sun we sat under the boat's hood.
\end{quote}

The curved roof of bamboo and matting over the middle of a river boat.

\begin{cousinweb}
\textbf{\bn{ছই}} (\emph{chhôi}) is inherited: CDIAL lists Bengali \emph{chai}, \textquotedblleft{}the thatched covering of a boat or cart\textquotedblright{}, under Sanskrit \emph{chadís}, \textquotedblleft{}a cover, a roof\textquotedblright{} (CDIAL 4980).
\end{cousinweb}

\begin{grammarlens}[title={What you've built}]
One more word for getting around.
\end{grammarlens}

\subsection*{Your turn}
\begin{itemize}
\raggedright
  \item \emph{Say it:} \emph{chhôi}
  \item \emph{Say it:} \emph{chhôi}, once more
  \item \emph{Say it:} \emph{rode āmrā noukār chhôiyer niche bôslām.} — In the sun we sat under the boat's hood.
  \item \emph{Recall:} say the Bengali for an oar, then the Bengali for an anchor, then say \emph{chhôi} again
\end{itemize}

\begin{tcolorbox}[breakable,colback=teal!4,colframe=teal!35!black,title={Before you move on}]
What does \bn{ছই} mean? (\textquotedblleft{}The hood of a country boat\textquotedblright{}.) Say it once more.
\end{tcolorbox}
\section[\texorpdfstring{\bn{বাওয়া} (bāoyā)}{বাওয়া (bāoyā)}]{\bn{বাওয়া} (bāoyā) — to row; to steer (a boat)}

\label{lesson:BN-C396-baoya}



\begin{quote}
Before the new one: say the Bengali for an anchor, then the Bengali for the hood of a country boat.
\end{quote}

\subsection*{You'll want to know: \bn{বাওয়া}}
\textbf{\bn{বাওয়া}} — \emph{bāoyā} — \textquotedblleft{}to row; to steer (a boat)\textquotedblright{}.

\begin{quote}
\bn{মাঝি} \bn{সারারাত} \bn{নৌকা} \bn{বাইল।}
\end{quote}

\begin{quote}
\emph{mājhi sārārāt noukā bāilô.}
\end{quote}

\begin{quote}
The boatman rowed the boat all night.
\end{quote}

Past (he/she …): \bn{তিনি} \bn{নৌকা} \bn{বাইলেন}, \emph{tini noukā bāilen}.

The stem is \bn{বা}-: \bn{আমি} \bn{বাই}, \bn{সে} \bn{বায়}, \bn{সে} \bn{বাইল}.

\begin{cousinweb}
\textbf{\bn{বাওয়া}} (\emph{bāoyā}) is inherited: CDIAL lists Bengali \emph{baoyā} and \emph{bahā}, \textquotedblleft{}to carry, to row, to steer\textquotedblright{}, under Sanskrit \emph{váhati}, \textquotedblleft{}carries, conveys\textquotedblright{} (CDIAL 11453).
\end{cousinweb}

\begin{grammarlens}[title={What you've built}]
One more word for getting around.
\end{grammarlens}

\subsection*{Your turn}
\begin{itemize}
\raggedright
  \item \emph{Say it:} \emph{bāoyā}
  \item \emph{Say it:} \emph{bāoyā}, once more
  \item \emph{Say it:} \emph{mājhi sārārāt noukā bāilô.} — The boatman rowed the boat all night.
  \item \emph{Recall:} say the Bengali for an anchor, then the Bengali for the hood of a country boat, then say \emph{bāoyā} again
\end{itemize}

\begin{tcolorbox}[breakable,colback=teal!4,colframe=teal!35!black,title={Before you move on}]
What does \bn{বাওয়া} mean? (\textquotedblleft{}To row; to steer (a boat)\textquotedblright{}.) Say it once more.
\end{tcolorbox}
